\section{Six Challenges}
\label{sec:1.3}

To understand the design and motivation behind different Monte Carlo methods, it is essential
to first appreciate the main challenges they aim to address. This section provides a high-level
overview of some of these challenges. Clearly, no single algorithm can effectively address all
of them simultaneously; otherwise, there would be no need for the many types of Monte Carlo
methods presented in this textbook! Throughout, each new algorithm we introduce will be accompanied
by a discussion of its pros and cons in relation to addressing different challenges, thus helping
to clarify its role within the vast landscape of Monte Carlo methods.

\paragraph{Challenge 1: High Dimensions}
Monte Carlo integration methods remain effective even in high-dimensional settings, provided that
one can obtain a sample from the target distribution. Unfortunately, sampling from a generic target
distribution in high dimensions is often difficult. On the bright side, there are important classes
of target distributions for which scalable sampling algorithms are available. For instance,
log-concave distributions can be sampled efficiently using algorithms closely related to those
developed for convex optimization (Chapter~6). As another example, if one can find a
\textit{proposal distribution} that is both close to the target and easy to sample from, then
samples from the proposal can be transformed into samples from the target (Chapters~2 and~3).

\paragraph{Challenge 2: Evaluating the Target and its Gradient}
The guiding applications to Bayesian statistics and statistical mechanics in Section~\ref{sec:1.2}
illustrate that the target distribution is often known only up to a normalizing constant. In many
Bayesian inference problems, evaluating the likelihood function can also be challenging, which has
motivated the development of approximate Bayesian computation, likelihood-free methods, and Monte
Carlo algorithms tailored to large datasets (Chapter~2). Relatedly, some Monte Carlo methods
leverage the gradient of the logarithm of the target density to improve the scalability to high
dimensions (Chapters~6 and~8), but evaluating these gradients can be computationally expensive.

\paragraph{Challenge 3: Multi-Modality}
Many sampling algorithms studied in this textbook are \textit{local}, in the sense that each new
draw typically lies in a small neighborhood of the previous one (Chapters~4, 6, and~8). For these
methods, sampling multi-modal target distributions poses a significant challenge, especially when
regions of high target density are separated by regions of low density. In this setting, local
methods may over-sample around one mode without adequately exploring others. A variety of strategies
have been developed to address this issue. For example, one can introduce flattened or \textit{tempered}
ancillary target distributions to accelerate exploration (Chapter~7), or augment the state space
with auxiliary variables to define an extended target in a higher-dimensional space with a more
benign landscape (Chapters~7 and~8).

\paragraph{Challenge 4: Computing Rare Events}
A different set of challenges arises when the test function $h(x) = \mathbf{1}_A(x)$ is the
indicator function of a small probability event $A$. In this case, a random sample from the target
distribution will typically contain no draws in $A$, causing the Monte Carlo estimate of $\Prob(A)$
to be zero. To obtain an estimate with low \textit{relative error}, it is then essential to sample
from a suitable proposal distribution rather than the target (Chapters~3 and~9). How should one
select the proposal distribution to minimize the variance of Monte Carlo estimates of rare events?
You will find out in Chapter~3.

\paragraph{Challenge 5: Choice of Variables and Conditioning of the Target}
For some algorithms, a key challenge is sampling joint distributions with strongly correlated
variables (Chapter~5) or variables with different scales (Chapters~4, 6, and~8). To mitigate
these issues, it is often beneficial to reparameterize the problem before applying Monte Carlo
methods. Other strategies include blocking together correlated variables, appropriately
preconditioning sampling algorithms, and, relatedly, using affine-invariant methods.

\paragraph{Challenge 6: Assessing Convergence}
Assessing the convergence of Monte Carlo methods can be challenging. Diagnostics based on
hypothesis tests can help rule out convergence, but cannot guarantee convergence (Chapter~4).
This challenge is particularly acute for multi-modal target distributions, where an algorithm
may appear to have converged after thoroughly exploring one mode while failing to visit others.
